\section{Trajectoire électronique, spectre de retour et conservation de la masse
}
\label{vi:sec:trayectoria-espectral-masa}

Cette section développe, sur l’opérateur de masse déjà construit, la partie
dynamique qui conserve la trajectoire électronique. Elle n’introduit ni seconde
définition de l’électron ni calibration supplémentaire. La chaîne causale est la
même que celle de l’article :

\[
 \mathrm{APP}\longrightarrow\mathrm{TRIT}\longrightarrow\mathrm{TPK}
 \longrightarrow X_{\rm enr}
 \longrightarrow\mathcal C_{\rm cont}^{\rm disc}
 \longrightarrow (\Gamma_e,\tau_e,\sigma_e)
 \longrightarrow M^\beta .
\]
Les coordonnées $\pi_{\rm HMT},\varphi_{\rm HMT},e_{\rm HMT}$ et
$\alpha_{\rm HMT}$ sont des sorties HMT corrélées engendrées à partir de l’état
enrichi ; elles ne sont pas introduites comme données conventionnelles. L’analyse de Fourier
qui apparaît plus bas s’applique au mot déjà produit et ne sélectionne pas la route.


\subsection{Objets centraux distincts et mémoire de la fibre
}

La cellule fixe $(5,5)$, la donnée de frontière $111111$, l’émission transversale
$555555$ et l’orientation $\tau_e=(+++---)^2$ ne sont pas des noms d’un même
objet. L’involution cellulaire

\[
 \rho_c(r,c)=(10-r,10-c)
\]
possède un unique point fixe, $(5,5)$. En celui-ci,

\[
 5+5=1+9,\qquad 5\cdot5=7+2\cdot9,
\]
tandis qu’en $(8,2)$ et $(2,8)$ la somme est conservée mais le quotient
du produit change. La cellule fixe détermine l’exception centrale ; le registre de
frontière et l’émission transversale conservent d’autres coordonnées de l’état.


En particulier, la fibre de $555555$ dans le canal de produit contient
$24$ positions orientées. Après une avancée du TPK, elle se distribue en

\[
 555177,\qquad 555717,\qquad 555771,
\]
avec huit positions par sortie. La valeur suivante n’est pas fonction de la signature
scalaire isolée : pour reconstruire l’évolution, il faut conserver la classe de
mémoire et l’orientation. Cette distinction est la raison pour laquelle l’opérateur
de masse agit sur des fibres de routes et non sur une liste de chiffres.


\subsection{Support cellulaire exact de la trajectoire centrale
}

La trajectoire de~\eqref{vi:dep:trayectoria-electron} commence en
$(r,c,h,v)=(5,5,1,1)$. Chaque bloc TRIT $(+1,-1,0)$ fait avancer une fois chaque
coordonnée et conserve un pas de seuil. Aux extrémités de chaque bloc, on
a $r=c$ ; au pas intermédiaire,

\[
 c-r\equiv h\pmod 9.
\]
Neuf blocs parcourent toutes les positions de la diagonale correspondante.
Les inversions simultanées de $(h,v)$ aux pas $27,54,81$ ajoutent la
diagonale conjuguée sans réinitialiser l’état. Par conséquent, la projection
cellulaire des $108$ pas a pour support exact

\begin{equation}
 \operatorname{Supp}_{\rm cell}(\Gamma_e)
 =\left\{(r,c)\in\{1,\ldots,9\}^2:
 c-r\equiv0,1,-1\pmod9\right\},
 \qquad
 \left|\operatorname{Supp}_{\rm cell}(\Gamma_e)\right|=27.
 \label{vi:eq:soporte-electron-27}
\end{equation}
L’égalité s’obtient par la règle de mise à jour, non par un comptage
rétrospectif. Elle décrit une morphologie discrète de trajectoire dans le tore
positionnel. Sa réalisation comme ruban de Möbius requiert en outre le transport
d’orientation et le recollement de la section~\ref{vi:sec:conjugacion} ; on
n’identifie pas le support fini à une forme matérielle tridimensionnelle.


\subsection{Deux histoires, une orientation sectorielle et une reconstruction exacte
}

La biographie centrale contient les histoires

\begin{equation}
 e_{++}=(2,2,5,-4,-3,-2)^2,
 \qquad
 e_{--}=(4,3,2,-2,-2,-5)^2.
 \label{vi:eq:historias-electronicas}
\end{equation}
Toutes deux s’expriment dans la coordinatisation sectorielle fixe
$\Sigma_{12}=(+++---)^2$. L’orientation absolue de la récurrence se
transporte séparément ; on n’identifie pas le signe d’une histoire à
l’étiquette sectorielle. Les deux histoires ont une somme nulle et le même histogramme
de valeurs absolues, mais pas la même mémoire positionnelle.


Pour une histoire $e=(e_0,\ldots,e_{11})$, nous définissons

\begin{equation}
 p_i=e_i+e_{i+6},\qquad a_i=e_i-e_{i+6},
 \qquad s_C=\sum_{i=0}^{5}p_i,
 \qquad M_i=p_i-p_5\quad(0\leq i<5).
 \label{vi:eq:codificador-memoria-electron}
\end{equation}
Le décodeur est

\begin{equation}
 p_5=\frac{s_C-\sum_{i=0}^{4}M_i}{6},\quad
 p_i=p_5+M_i,\quad
 e_i=\frac{p_i+a_i}{2},\quad e_{i+6}=\frac{p_i-a_i}{2}.
 \label{vi:eq:decodificador-memoria-electron}
\end{equation}
Les conditions d’intégralité sont la divisibilité par six de la première
expression et la parité commune de $p_i$ et de $a_i$. Pour
\eqref{vi:eq:historias-electronicas}, on obtient

\[
 M^{++}=(8,8,14,-4,-2),\qquad
 M^{--}=(18,16,14,6,6),\qquad a=0.
\]
Substituer ces données dans~\eqref{vi:eq:decodificador-memoria-electron}
restitue les douze composantes. La distinction des histoires est ainsi prouvée
même si une lecture partielle coïncide. On n’en conclut pas que leurs masses doivent être
différentes : un lecteur de masse peut être invariant sous un transport qui
conserve moins de coordonnées que l’état enrichi.


\subsection{Identité spectrale du lecteur de retour
}

Soit $\tau=(\tau_0,\ldots,\tau_{11})\in\{-1,+1\}^{12}$ et soient les indices
cycliques. Nous définissons

\begin{equation}
 C_s(\tau)=\sum_{j=0}^{11}\tau_j\tau_{j+s},
 \qquad
 \Omega(\tau)=-2C_1-4C_2-6C_3,
 \qquad
 P_6(\tau)=\frac{C_6}{2}.
 \label{vi:eq:lectores-correlacion}
\end{equation}
La première égalité est la forme développée de
$\Omega=4A_3-3A_4$ déjà démontrée dans
\eqref{vi:dep:KOmega}. Pour rendre visibles les modes qui transportent la
correction, nous introduisons, après avoir produit $\tau$,

\begin{equation}
 T_k(\tau)=\sum_{j=0}^{11}\tau_j
 e^{-2\pi i k j/12},\qquad 0\leq k<12.
 \label{vi:eq:dft-retorno}
\end{equation}

\begin{proposition}[Décomposition spectrale exacte
]
Pour tout mot réel de longueur douze,

\begin{align}
 C_s&=\frac1{12}\sum_{k=0}^{11}|T_k|^2
          \cos\!\left(\frac{2\pi ks}{12}\right),\label{vi:eq:wiener-khinchin-finito}\\
 \Omega&=\frac1{12}\sum_{k=0}^{11}|T_k|^2w_k,
 \quad
 w_k=-2\cos\theta_k-4\cos2\theta_k-6\cos3\theta_k,
 \quad\theta_k=\frac{2\pi k}{12},\label{vi:eq:omega-espectral}\\
 P_6&=\frac1{24}\sum_{k=0}^{11}|T_k|^2(-1)^k.
 \label{vi:eq:p6-espectral}
\end{align}
\end{proposition}
\begin{proof}
En développant $|T_k|^2$ apparaît
$\sum_{j,\ell}\tau_j\tau_\ell e^{-2\pi ik(j-\ell)/12}$.
La somme sur $k$ vaut douze lorsque $j-\ell$ coïncide avec le déplacement
prescrit et zéro sinon. C’est l’orthogonalité finie des caractères
et elle démontre~\eqref{vi:eq:wiener-khinchin-finito}. Insérer successivement
$s=1,2,3$ dans~\eqref{vi:eq:lectores-correlacion} produit
\eqref{vi:eq:omega-espectral} ; $s=6$ produit
\eqref{vi:eq:p6-espectral}.

\end{proof}

Pour $\tau_e=(+++---)^2$, on peut écrire, avec
$z_k=e^{-2\pi ik/12}$,

\begin{equation}
 T_k=(1+(-1)^k)(1-z_k^3)(1+z_k+z_k^2).
 \label{vi:eq:factorizacion-tk-electron}
\end{equation}
Seuls $k=2,6,10$ subsistent, avec

\[
 T_2=4-4i\sqrt3,\qquad T_6=4,\qquad
 T_{10}=4+4i\sqrt3,
\]
et, par conséquent,
\begin{equation}
 |T_2|^2=64,\qquad |T_6|^2=16,\qquad |T_{10}|^2=64.
 \label{vi:eq:potencias-espectrales-electron}
\end{equation}
Les poids $w_2,w_6,w_{10}$ sont $7,4,7$. En conséquence,

\begin{equation}
 \Omega(\tau_e)=\frac{64\cdot7+16\cdot4+64\cdot7}{12}=80,
 \qquad P_6(\tau_e)=6.
 \label{vi:eq:omega-p6-electron}
\end{equation}
Cela récupère le second lecteur de~\eqref{vi:dep:KOmega} à partir des modes
de la trajectoire. Le mot alterné $(+-)^6$, qui conserve le même
histogramme de signes, a $|T_6|^2=144$ et produit $\Omega=48$ ; il constitue
un contrôle négatif du fait que l’ordre, et pas seulement le comptage, intervient.
Les translations cycliques multiplient $T_k$ par une phase et conservent
$|T_k|^2$, de sorte que le lecteur ne récupère pas à lui seul l’origine temporelle.


La contribution de retour à l’exposant électronique s’exprime par

\begin{equation}
 \kappa_e=
 \frac1{12}\sum_{k=0}^{11}|T_k|^2
 \left(w_k+\frac{\Delta_4}{2}(-1)^k\right)-54\alpha_{\rm HMT}
 =80-54\alpha_{\rm HMT}+6\Delta_4.
 \label{vi:eq:kappa-espectral-electron}
\end{equation}
Cette identité n’introduit pas d’hamiltonien spatial : elle démontre que le
correcteur de masse déjà défini est une fonctionnelle spectrale exacte du mot
de retour.


\subsection{Extension de l’équation électronique aux multisections
}

Le registre d’une route $\gamma$ fournit
$\sigma_\gamma=(n_A,s_C,k_\Delta,b_{120},b_\zeta)$ et les coordonnées de
tour $\eta_\gamma=(\eta_h)_h$. Dans une coordinatisation commune, le caractère raffiné est

\begin{equation}
 \begin{aligned}
 \log\chi_{\rm ref}(\sigma_\gamma,\eta_\gamma)
 &=n_Ax+\frac{s_C}{6}y+k_\Delta\Delta_4+b_\zeta\zeta_{270}
   +b_{120}s_{120}+\sum_{h\in\langle90,120\rangle}\eta_hs_h\\
 &=n_Ax+\frac{s_C}{6}y+k_\Delta\Delta_4+b_\zeta\zeta_{270}
   +\sum_{h\in\langle90,120\rangle}N_hs_h,\\
 N_h&:=\eta_h+b_{120}\mathbf1_{\{h=120\}}.
 \end{aligned}
 \label{vi:eq:caracter-refinado-explicito}
\end{equation}
où $x=A_{\rm rad}$, $y=C^*_{\rm rad}$,
$\zeta_{270}=135\Delta_4^2$ et $s_h=q_-^h-q_+^h$. Les deux lignes sont la
même évaluation : $b_{120}$ n’apparaît qu’une seule fois, comme terme explicite ou
absorbé dans $N_{120}$. Pour une tour infinie, on exige la convergence absolue
de la dernière série. Cette formule conserve séparément la torsion $\zeta_{270}$,
les deux feuillets et les deux compositions de hauteur $360=4\cdot90=3\cdot120$.


Dans la base de routes d’une fibre, soit
$B^r e_\gamma=\kappa_\gamma^r e_\gamma$, avec $r=D$ ou $\Omega$. La loi
opératorielle de la section~\ref{vi:sec:operador-masa} donne, dans sa coordinatisation diagonale,

\begin{equation}
 \boxed{
 \frac{m_\gamma^\beta}{m_e^\beta}=
 \chi_{\rm ref}(\sigma_\gamma-\sigma_e,
                 \eta_\gamma-\eta_e)
 R_{\rm act}^{\kappa_\gamma-\kappa_e}.}
 \label{vi:eq:ley-masas-trayectoria}
\end{equation}
L’équation électronique est la section de référence de cette identité, non
un chiffre expérimental utilisé pour choisir des routes. Les étiquettes physiques sont
attribuées après la construction des fibres, des représentations et des canaux ; la ressemblance
numérique ne remplace pas cet appariement prospectif.


\subsection{Critère opératoriel de transport à masse fixe
}

Nous définissons la coordonnée adimensionnelle

\begin{equation}
 \Lambda_\gamma:=\log\frac{m_\gamma^\beta}{m_e^\beta}
 \label{vi:eq:lambda-masa-ruta}
\end{equation}
dans une coordinatisation commune avec des masses positives. Ce symbole n’est ni la longueur
réciproque $\Lambda=L_\alpha X^{-1}$ ni la coordonnée
$\bar\lambda_\Gamma$ employée dans d’autres sections.


Soit $\mathscr D$ un domaine fini de routes stable sous un transport TPK
bijectif $F:\mathscr D\to\mathscr D$, et soit
$V_Fe_\gamma=e_{F\gamma}$ son opérateur de permutation. Alors

\begin{equation}
 [M^\beta,V_F]e_\gamma
 =m_e^\beta\left(e^{\Lambda_{F\gamma}}-e^{\Lambda_\gamma}\right)
 e_{F\gamma}.
 \label{vi:eq:conmutador-masa-transporte}
\end{equation}
Par conséquent,
\begin{equation}
 [M^\beta,V_F]=0
 \quad\Longleftrightarrow\quad
 \Lambda_{F\gamma}=\Lambda_\gamma
 \quad\text{pour toute }\gamma\in\mathscr D.
 \label{vi:eq:criterio-conservacion-masa}
\end{equation}
En utilisant~\eqref{vi:eq:caracter-refinado-explicito}, la condition est

\begin{equation}
 \delta n_Ax+\frac{\delta s_C}{6}y
 +\delta k_\Delta\Delta_4+\delta b_\zeta\zeta_{270}
 +\sum_h\delta N_hs_h
 +\delta\kappa\log R_{\rm act}=0.
 \label{vi:eq:balance-conservacion-masa}
\end{equation}
Ici,
\(
 \delta N_h=\delta\eta_h+
 \delta b_{120}\mathbf1_{\{h=120\}},
\)
de sorte que la différence conserve les coordonnées de tour originales et
ajoute le terme d’incidence de $b_{120}$ exactement une fois.
Elle se vérifie après avoir produit $F$ au moyen du TPK ; on ne sélectionne pas
$F$ rétrospectivement par la valeur de la masse. Le critère permet que
les événements et la mémoire changent pourvu que la coordonnée totale soit conservée.
Il n’affirme ni que tout transport TPK conserve la masse ni qu’il épuise la dynamique
d’un système lié. Ses hypothèses sont : domaine stable, transport
bijectif, coordinatisation fixe, masses positives et convergence des tours.


\subsection{Résultat conjoint
}

Les équations~\eqref{vi:eq:soporte-electron-27},
\eqref{vi:eq:decodificador-memoria-electron},
\eqref{vi:eq:kappa-espectral-electron},
\eqref{vi:eq:ley-masas-trayectoria} et
\eqref{vi:eq:criterio-conservacion-masa} ferment une même chaîne. La
trajectoire centrale parcourt exactement vingt-sept cellules ; ses histoires
conservent une mémoire réversible ; le retour possède trois modes non nuls qui
produisent $(\Omega,P_6)=(80,6)$ ; la loi des masses étend la section
électronique aux multisections ; et le commutateur identifie exactement
quels transports préservent la masse. Morphologie cellulaire, spectre de retour et
grandeur sont des lecteurs coordonnés du même état enrichi, non des
observables interchangeables.

